\subsection{显著单词的刻画}

一个单词 \(A \in \mathbf{L}_0(\mathbf{S})\) 如果它满足以下两个性质，则称其为平衡的：

\begin{enumerate}
\def\labelenumi{(\arabic{enumi})}
\item
  \(l(A) = n(A) + 1\) （这意味着 \(A\) 非空）。
\item
  对于每一个真初始段 \(B\) （即 \(A\) 的真初始段）, \(l(B) \leqslant n(B)\) .
\end{enumerate}

命题2. 一个词是显著的当且仅当它是平衡的。

设 \(A\) 是一个属于显著序列的显著词，

\[
A _ {1}, A _ {2}, \dots , A _ {n}.
\]

我们将通过对 k 进行归纳来证明每个 \(A_{k}\) 是平衡的。假设这一点已对 \(A_{j}\) （指标 j \textless{} k）成立，现在让我们证明它对 \(A_{k}\) 成立。如果 \(A_{k}\) 是权重为 0 的符号（当 k = 1 时这是唯一可能的情况）， \(A_{k}\) 是平衡的，因为

\[
l (A _ {k}) = 1 \quad \text { 且 } \quad n (A _ {k}) = 0.
\]

如果 \(A_{k}\) 不是权为0的符号，那么 \(A_{k} = fB_{1}B_{2}\ldots B_{p}\) ，其中 \(f\) 是权重为 \(p\) 的符号，并且 \(B_{j}\) 形如 \(A_{i_{j}}\) ，其中 \(i_{j} < k\) ，因此根据归纳假设它们是平衡词。我们有

\[
\begin{array}{r l} l (A _ {k}) & = 1 + l (B _ {1}) + l (B _ {2}) + \dots + l (B _ {p}) \\ & = 1 + (n (B _ {1}) + 1) + (n (B _ {2}) + 1) + \dots + (n (B _ {p}) + 1) \\ & = 1 + p + n (B _ {1}) + n (B _ {2}) + \dots + n (B _ {p}) \\ & = 1 + n (A _ {k}). \end{array}
\]

设 \(C\) 是 \(A\) 的一个真初始段，并且设 \(q\) 为整数 \(m < p\) 中最大的一个，使得 \(B_{m}\) 是 \(C\) 的一个段，因此

\[
C = f B _ {1} B _ {2} \dots B _ {q} D,
\]

其中 D 是 \(B_{q+1}\) 的一个真初始段。然后

\[
\begin{array}{r l} & l (C) = 1 + l (B _ {1}) + l (B _ {2}) + \dots + l (B _ {q}) + l (D) \\ & \leqslant 1 + (n (B _ {1}) + 1) + (n (B _ {2}) + 1) + \dots + (n (B _ {q}) + 1) + n (D) \\ & \leqslant p + n (B _ {1}) + \dots + n (B _ {q}) + n (D) = n (C). \end{array}
\]

因此， \(A_{k}\) 是平衡的。

为了证明反过来每个平衡词都是显著的，我们需要以下两个引理：

引理1. 设 \(A\) 为一个平衡词。则对每个整数 \(k\) ，只要 \(0 \leqslant k < l(A)\) ，就恰好存在一个平衡段 \(S\) （ \(A\) 的段），它始于第 \((k + 1)\) 位。

S的唯一性是以下论断的直接推论：若T是一个平衡词，则根据定义，T的任何真初始段都不是平衡的。下面证明S的存在性。将A写作A = BC，其中 \(l(B) = k\) 。对于每个满足 \(0 \leqslant i \leqslant q = l(C)\) 的 i ，令 \(C_i\) 为C的长度为i的初始段。由于B是A的真初始段，我们有

\[
l (C _ {q}) = l (A) - l (B) \geqslant n (A) + 1 - n (B) = n (C _ {q}) + 1.
\]

另一方面，我们有 \(0 = l(C_0) \leqslant n(C_0) = 0\) . 设 \(i\) 为最大的整数 \(j < q\) ，使得 \(l(C_h) \leqslant n(C_h)\) 对 \(0 \leqslant h \leqslant j\) 成立；于是我们有 \(l(C_i) \leqslant n(C_i)\) 并且 \(l(C_{i+1}) \geqslant n(C_{i+1}) + 1\) 。我们证明 \(C_{i+1}\) 是平衡的：与真初始段相关的条件因 \(i\) 的定义而得到满足；另一方面，我们有

\[
n (C _ {i + 1}) + 1 \leqslant l (C _ {i + 1}) = l (C _ {i}) + 1 \leqslant n (C _ {i}) + 1 \leqslant n (C _ {i + 1}) + 1,
\]

使得 \(l(C_{i+1}) = n(C_{i+1}) + 1\) ，引理1的证明完成。

引理2. 每个平衡词A都可以写成如下形式

\[
A = f A _ {1} A _ {2} \dots A _ {p},
\]

其中 \(A_{i}\) 是平衡的，且 \(n(f) = p\) .

设 \(f\) 是 \(A\) 的初始符号。由引理1， \(A\) 可以写成

\[
f A _ {1} A _ {2} \dots A _ {p},
\]

其中 \(A_{i}\) 是平衡的：我们只需归纳地定义 \(A_{i}\) 为从第 \(k(i)\) 位开始的A的平衡段，其中

\[
k (i) = 2 + \sum_ {j <   i} l (A _ {j}).
\]

此外，我们有

\[
\begin{array}{r l} 1 + l (A _ {1}) + \dots + l (A _ {p}) & = l (A) = n (A) + 1 \\ & = n (f) + n (A _ {1}) + \dots + n (A _ {p}) + 1 \\ & = n (f) + (l (A _ {1}) - 1) + \dots + (l (A _ {p}) - 1) + 1, \end{array}
\]

由此可得 \(n(f) = p\) .

既然两个引理已经证明，那么通过对A的长度进行归纳，显然每个平衡词A都是显著的，这由引理2和命题1可知。

推论1. 设 \(A\) 是一个显著的词。对于每个整数 \(k\) ，只要 \(0 \leqslant k < l(A)\) ，就恰好存在 \(A\) 的一个显著的段，它始于第 \((k + 1)\) 位。

推论2. 每个显著的词都可以以唯一的方式写成形式 \(fA_{1}A_{2}\ldots A_{p}\) ，其中 \(A_{i}\) 是显著的，且 \(n(f) = p\).*

